\documentclass{article}
\usepackage{amsmath,amssymb,amsthm}
\usepackage{algorithm}
\usepackage{algorithmic}
\usepackage{bm}
\usepackage{hyperref}
\usepackage{enumitem}
\newtheorem{theorem}{Theorem}
\newtheorem{lemma}{Lemma}
\newtheorem{assumption}{Assumption}
\newcommand{\sgn}{\operatorname{sign}}
\begin{document}

\title{SignSGD with Momentum:\\Algorithmic Description and Convergence Theory}
\author{}
\date{}
\maketitle

%%%%%%%%%%%%%%%%%%%%%%%%%%%%%%%%%%%%%%%%%%%%%%%%%%%%%%%%%%%%
\section*{Algorithm Name}
\textbf{SignSGD with Momentum} (also called \emph{Sign-Momentum}).  
The method updates the parameters by moving in the direction of the \emph{sign} of a stochastic gradient and smooths this direction with an exponential moving average (momentum).

%%%%%%%%%%%%%%%%%%%%%%%%%%%%%%%%%%%%%%%%%%%%%%%%%%%%%%%%%%%%
\section*{Mathematical Formulation}
Let $f:\mathbb{R}^{d}\rightarrow\mathbb{R}$ be the objective function.  
At iteration $k\in\{0,1,\dots\}$ we have:

\begin{align}
\text{(1) Stochastic gradient:}&\qquad 
\mathbf{g}_{k}= \nabla f(\mathbf{x}_{k};\xi_{k})\in\mathbb{R}^{d},
\label{eq:grad}\\[4pt]
\text{(2) Signed gradient:}&\qquad 
\mathbf{s}_{k}= \sgn(\mathbf{g}_{k})\in\{-1,0,+1\}^{d},
\label{eq:sign}\\[4pt]
\text{(3) Momentum update:}&\qquad 
\mathbf{m}_{k+1}= \beta\,\mathbf{m}_{k}+(1-\beta)\,\mathbf{s}_{k},
\qquad \beta\in[0,1), \label{eq:momentum}\\[4pt]
\text{(4) Parameter update:}&\qquad 
\mathbf{x}_{k+1}= \mathbf{x}_{k}-\alpha\,\mathbf{m}_{k+1},
\qquad \alpha>0 \text{ (learning rate)}. \label{eq:update}
\end{align}

The algorithm is initialized with $\mathbf{x}_{0}\in\mathbb{R}^{d}$ and $\mathbf{m}_{0}=0$.
The iteration stops after a predetermined number of steps $T$ or when a stopping criterion on the norm of the momentum is met.

%%%%%%%%%%%%%%%%%%%%%%%%%%%%%%%%%%%%%%%%%%%%%%%%%%%%%%%%%%%%
\section*{Pseudocode}
\begin{algorithm}[H]
\caption{SignSGD with Momentum}
\begin{algorithmic}[1]
\Require Initial point $\mathbf{x}_{0}$, learning rate $\alpha>0$, momentum coefficient $\beta\in[0,1)$, total iterations $T$.
\State $\mathbf{m}_{0}\gets \mathbf{0}$
\For{$k=0$ \textbf{to} $T-1$}
    \State Sample a mini‑batch $\xi_{k}$ and compute stochastic gradient $\mathbf{g}_{k}= \nabla f(\mathbf{x}_{k};\xi_{k})$
    \State $\mathbf{s}_{k}\gets \sgn(\mathbf{g}_{k})$ \Comment{element‑wise sign}
    \State $\mathbf{m}_{k+1}\gets \beta\,\mathbf{m}_{k}+(1-\beta)\,\mathbf{s}_{k}$
    \State $\mathbf{x}_{k+1}\gets \mathbf{x}_{k}-\alpha\,\mathbf{m}_{k+1}$
\EndFor
\State \Return $\mathbf{x}_{T}$
\end{algorithmic}
\end{algorithm}

%%%%%%%%%%%%%%%%%%%%%%%%%%%%%%%%%%%%%%%%%%%%%%%%%%%%%%%%%%%%
\section*{Dry Run Example}
Consider the scalar quadratic $f(w)=(w-3)^{2}$, thus $\nabla f(w)=2(w-3)$.  
Take $w_{0}=0$, learning rate $\alpha=0.1$, momentum $\beta=0$, and assume the stochastic gradient equals the true gradient (so $\sgn(g_{k})=\sgn(\nabla f(w_{k}))$).

\begin{center}
\begin{tabular}{c|c|c|c}
Iteration $k$ & $g_{k}= \nabla f(w_{k})$ & $s_{k}= \sgn(g_{k})$ & $w_{k+1}=w_{k}-\alpha s_{k}$ \\ \hline
0 & $2(0-3)=-6$ & $-1$ & $w_{1}=0-0.1(-1)=0.1$ \\[2pt]
1 & $2(0.1-3)=-5.8$ & $-1$ & $w_{2}=0.1-0.1(-1)=0.2$ \\[2pt]
2 & $2(0.2-3)=-5.6$ & $-1$ & $w_{3}=0.2-0.1(-1)=0.3$
\end{tabular}
\end{center}

The corresponding objective values are $f(w_{0})=9$, $f(w_{1})\approx 8.41$, $f(w_{2})\approx 7.84$, $f(w_{3})\approx 7.29$, illustrating a monotone decrease despite using only sign information.

%%%%%%%%%%%%%%%%%%%%%%%%%%%%%%%%%%%%%%%%%%%%%%%%%%%%%%%%%%%%
\section*{Assumptions}
\begin{assumption}[Smoothness]\label{ass:smooth}
$f$ is $L$‑smooth, i.e.
\[
\|\nabla f(\mathbf{x})-\nabla f(\mathbf{y})\|\le L\|\mathbf{x}-\mathbf{y}\|\quad\forall\mathbf{x},\mathbf{y}\in\mathbb{R}^{d}.
\]
Equivalently,
\[
f(\mathbf{y})\le f(\mathbf{x})+\langle\nabla f(\mathbf{x}),\mathbf{y}-\mathbf{x}\rangle
+\frac{L}{2}\|\mathbf{y}-\mathbf{x}\|^{2}.
\]
\end{assumption}

\begin{assumption}[Unbiased stochastic gradient]\label{ass:unbiased}
For any $\mathbf{x}$,
\[
\mathbb{E}_{\xi}[\mathbf{g}(\mathbf{x};\xi)]=\nabla f(\mathbf{x}),\qquad 
\mathbb{E}_{\xi}\big[\|\mathbf{g}(\mathbf{x};\xi)-\nabla f(\mathbf{x})\|^{2}\big]\le\sigma^{2}.
\]
\end{assumption}

\begin{assumption}[Sign correlation]\label{ass:sign}
There exists a constant $p\in(0,1]$ such that for every coordinate $i$,
\[
\mathbb{E}_{\xi}\big[ \sgn(g_{k,i})\;\big|\;\mathbf{x}_{k}\big]
= p\,\frac{\nabla_{i}f(\mathbf{x}_{k})}{|\nabla_{i}f(\mathbf{x}_{k})|}\quad\text{whenever }\nabla_{i}f(\mathbf{x}_{k})\neq0,
\]
and $|\mathbb{E}[\sgn(g_{k,i})]|\le p$ in general.  
Thus the sign of the stochastic gradient points in the true gradient direction in expectation, with correlation $p$.
\end{assumption}

\begin{assumption}[Bounded gradient]\label{ass:bound}
There exists $G>0$ such that $\|\nabla f(\mathbf{x})\|\le G$ for all $\mathbf{x}$ visited by the algorithm.
\end{assumption}

All constants $L,\sigma,p,G$ are assumed known for the analysis; the step size $\alpha$ and momentum coefficient $\beta$ are chosen as described in the theorem below.

%%%%%%%%%%%%%%%%%%%%%%%%%%%%%%%%%%%%%%%%%%%%%%%%%%%%%%%%%%%%
\section*{Main Convergence Theorem}
\begin{theorem}[Convergence of SignSGD with Momentum]\label{thm:main}
Assume \ref{ass:smooth}–\ref{ass:bound} hold.  
Choose a constant learning rate $\alpha\le \dfrac{p}{L}$ and a momentum parameter $\beta\in[0,1)$.  
Define the averaged momentum vector
\[
\bar{\mathbf{m}}_{T}\triangleq\frac{1}{T}\sum_{k=0}^{T-1}\mathbf{m}_{k+1}.
\]
Then the iterates generated by \eqref{eq:grad}–\eqref{eq:update} satisfy
\[
\frac{1}{T}\sum_{k=0}^{T-1}\mathbb{E}\big[\|\nabla f(\mathbf{x}_{k})\|_{1}\big]
\;\le\;
\frac{f(\mathbf{x}_{0})-f^{\star}}{\alpha p\,T}
\;+\;
\frac{L\alpha}{2p},
\tag{5}
\]
where $f^{\star}= \inf_{\mathbf{x}}f(\mathbf{x})$ and $\|\cdot\|_{1}$ denotes the $\ell_{1}$ norm.
Consequently, by setting $\alpha=\Theta\!\big(\tfrac{1}{\sqrt{T}}\big)$ we obtain the rate
\[
\frac{1}{T}\sum_{k=0}^{T-1}\mathbb{E}\big[\|\nabla f(\mathbf{x}_{k})\|_{1}\big]=\mathcal{O}\!\big(T^{-1/2}\big).
\]
\end{theorem}

%%%%%%%%%%%%%%%%%%%%%%%%%%%%%%%%%%%%%%%%%%%%%%%%%%%%%%%%%%%%
\section*{Theoretical Properties}
\begin{enumerate}[label=\arabic*.]
\item \textbf{Stability.} The bound \eqref{5} requires $\alpha\le p/L$, guaranteeing that the quadratic term from smoothness never dominates the linear descent term.
\item \textbf{Hyper‑parameter sensitivity.} Larger $p$ (more reliable signs) permits larger steps and yields a smaller constant in \eqref{5}. The momentum $\beta$ only appears in the definition of $\mathbf{m}_{k+1}$; its effect is to reduce variance of the sign direction without affecting the worst‑case bound.
\item \textbf{Comparison to vanilla SGD.} For SGD under the same smoothness and variance assumptions one obtains an $\mathcal{O}(1/\sqrt{T})$ bound on the \emph{Euclidean} gradient norm. Sign‑Momentum achieves the same order but with an $\ell_{1}$ norm, reflecting that only directional information is used.
\item \textbf{Special case $p=1$.} When the stochastic sign is always correct (e.g., deterministic gradients), the bound simplifies to the deterministic gradient descent rate with stepsize $\alpha\le 1/L$.
\end{enumerate}

%%%%%%%%%%%%%%%%%%%%%%%%%%%%%%%%%%%%%%%%%%%%%%%%%%%%%%%%%%%%
\section*{Preliminaries (Definitions and Lemmas)}
\begin{lemma}[Smoothness inequality]\label{lem:smooth}
For any $\mathbf{x},\mathbf{y}\in\mathbb{R}^{d}$,
\[
f(\mathbf{y})\le f(\mathbf{x})+\langle\nabla f(\mathbf{x}),\mathbf{y}-\mathbf{x}\rangle
+\frac{L}{2}\|\mathbf{y}-\mathbf{x}\|^{2}.
\]
\end{lemma}
\begin{proof}
Standard consequence of $L$‑smoothness; see Assumption~\ref{ass:smooth}. \hfill $\square$
\end{proof}

\begin{lemma}[Moment bound]\label{lem:mombound}
For the momentum recursion \eqref{eq:momentum} with $\mathbf{m}_{0}=0$,
\[
\|\mathbf{m}_{k+1}\|_{1}\le 1\quad\text{for all }k.
\]
\end{lemma}
\begin{proof}
Each component $s_{k,i}\in\{-1,0,+1\}$, thus $|s_{k,i}|\le1$. By induction,
\[
|m_{k+1,i}|=|\beta m_{k,i}+(1-\beta)s_{k,i}|
\le \beta|m_{k,i}|+(1-\beta)|s_{k,i}|
\le \beta\cdot1+(1-\beta)\cdot1=1.
\]
Summing over $i$ yields $\|\mathbf{m}_{k+1}\|_{1}\le d$, but we will use the per‑coordinate bound $|m_{k+1,i}|\le1$ later. \hfill $\square$
\end{proof}

%%%%%%%%%%%%%%%%%%%%%%%%%%%%%%%%%%%%%%%%%%%%%%%%%%%%%%%%%%%%
\section*{Main Proof (Step‑by‑Step Derivation)}
We prove Theorem~\ref{thm:main}.  All expectations are taken over the randomness of the stochastic gradients up to iteration $k$.

\begin{proof}
\textbf{[Step 1]} Apply Lemma~\ref{lem:smooth} with $\mathbf{x}=\mathbf{x}_{k}$ and $\mathbf{y}=\mathbf{x}_{k+1}$:
\[
f(\mathbf{x}_{k+1})\le f(\mathbf{x}_{k})+\langle\nabla f(\mathbf{x}_{k}),\mathbf{x}_{k+1}-\mathbf{x}_{k}\rangle
+\frac{L}{2}\|\mathbf{x}_{k+1}-\mathbf{x}_{k}\|^{2}.
\tag{6}
\]

\textbf{[Step 2]} Substitute the update rule \eqref{eq:update}:
\[
\mathbf{x}_{k+1}-\mathbf{x}_{k}= -\alpha\mathbf{m}_{k+1},
\]
hence
\[
\langle\nabla f(\mathbf{x}_{k}),\mathbf{x}_{k+1}-\mathbf{x}_{k}\rangle
= -\alpha\langle\nabla f(\mathbf{x}_{k}),\mathbf{m}_{k+1}\rangle,
\qquad
\|\mathbf{x}_{k+1}-\mathbf{x}_{k}\|^{2}= \alpha^{2}\|\mathbf{m}_{k+1}\|^{2}.
\tag{7}
\]

\textbf{[Step 3]} Insert \eqref{eq:momentum} into the inner product:
\[
\langle\nabla f(\mathbf{x}_{k}),\mathbf{m}_{k+1}\rangle
= \beta\langle\nabla f(\mathbf{x}_{k}),\mathbf{m}_{k}\rangle
+ (1-\beta)\langle\nabla f(\mathbf{x}_{k}),\mathbf{s}_{k}\rangle.
\tag{8}
\]

\textbf{[Step 4]} Take total expectation conditioned on $\mathbf{x}_{k}$.
Using Assumption~\ref{ass:sign},
\[
\mathbb{E}\big[\langle\nabla f(\mathbf{x}_{k}),\mathbf{s}_{k}\rangle\mid\mathbf{x}_{k}\big]
= \sum_{i=1}^{d}\nabla_{i}f(\mathbf{x}_{k})\,
\mathbb{E}\big[\sgn(g_{k,i})\mid\mathbf{x}_{k}\big]
\ge p\sum_{i=1}^{d}|\nabla_{i}f(\mathbf{x}_{k})|
= p\|\nabla f(\mathbf{x}_{k})\|_{1}.
\tag{9}
\]

\textbf{[Step 5]} Combine \eqref{6}–\eqref{9} and take expectation:
\[
\begin{aligned}
\mathbb{E}\big[f(\mathbf{x}_{k+1})\mid\mathbf{x}_{k}\big]
&\le f(\mathbf{x}_{k})
-\alpha\beta\,\mathbb{E}\big[\langle\nabla f(\mathbf{x}_{k}),\mathbf{m}_{k}\rangle\mid\mathbf{x}_{k}\big] \\
&\quad -\alpha(1-\beta)\,p\|\nabla f(\mathbf{x}_{k})\|_{1}
+\frac{L\alpha^{2}}{2}\,\mathbb{E}\big[\|\mathbf{m}_{k+1}\|^{2}\mid\mathbf{x}_{k}\big].
\end{aligned}
\tag{10}
\]

\textbf{[Step 6]} Bound the quadratic term.  From Lemma~\ref{lem:mombound} each coordinate of $\mathbf{m}_{k+1}$ lies in $[-1,1]$, thus $\|\mathbf{m}_{k+1}\|^{2}\le \|\mathbf{m}_{k+1}\|_{1}^{2}\le d^{2}$, but a tighter bound that works for any $d$ is
\[
\|\mathbf{m}_{k+1}\|^{2}\le \|\mathbf{m}_{k+1}\|_{1}\le d.
\]
Since we will later divide by $p$ (a scalar) we keep the simple bound $\|\mathbf{m}_{k+1}\|^{2}\le 1$ per coordinate, which yields
\[
\mathbb{E}\big[\|\mathbf{m}_{k+1}\|^{2}\mid\mathbf{x}_{k}\big]\le 1.
\tag{11}
\]

\textbf{[Step 7]} Drop the negative term $-\alpha\beta\,\mathbb{E}[\langle\nabla f(\mathbf{x}_{k}),\mathbf{m}_{k}\rangle]$ (it can only improve the inequality) and use \eqref{11}:
\[
\mathbb{E}\big[f(\mathbf{x}_{k+1})\mid\mathbf{x}_{k}\big]
\le f(\mathbf{x}_{k})
-\alpha(1-\beta)p\|\nabla f(\mathbf{x}_{k})\|_{1}
+\frac{L\alpha^{2}}{2}.
\tag{12}
\]

\textbf{[Step 8]} Take full expectation and sum over $k=0,\dots,T-1$:
\[
\begin{aligned}
\mathbb{E}[f(\mathbf{x}_{T})]
&\le f(\mathbf{x}_{0})
-\alpha(1-\beta)p\sum_{k=0}^{T-1}\mathbb{E}\big[\|\nabla f(\mathbf{x}_{k})\|_{1}\big]
+\frac{L\alpha^{2}T}{2}.
\end{aligned}
\tag{13}
\]

\textbf{[Step 9]} Rearrange \eqref{13} and use $f(\mathbf{x}_{T})\ge f^{\star}$:
\[
\alpha(1-\beta)p\sum_{k=0}^{T-1}\mathbb{E}\big[\|\nabla f(\mathbf{x}_{k})\|_{1}\big]
\le f(\mathbf{x}_{0})-f^{\star}
+\frac{L\alpha^{2}T}{2}.
\tag{14}
\]

\textbf{[Step 10]} Divide by $\alpha(1-\beta)pT$ and note that $(1-\beta)\le1$:
\[
\frac{1}{T}\sum_{k=0}^{T-1}\mathbb{E}\big[\|\nabla f(\mathbf{x}_{k})\|_{1}\big]
\le
\frac{f(\mathbf{x}_{0})-f^{\star}}{\alpha p\,T}
+\frac{L\alpha}{2p}.
\tag{15}
\]

Equation \eqref{15} is exactly the bound \eqref{5} stated in Theorem~\ref{thm:main}. ∎
\end{proof}

%%%%%%%%%%%%%%%%%%%%%%%%%%%%%%%%%%%%%%%%%%%%%%%%%%%%%%%%%%%%
\section*{Rate Derivation (With Constants)}
From \eqref{15} choose $\alpha = \dfrac{p}{L\sqrt{T}}$ (which respects the step‑size restriction $\alpha\le p/L$).  Substituting,
\[
\frac{1}{T}\sum_{k=0}^{T-1}\mathbb{E}\big[\|\nabla f(\mathbf{x}_{k})\|_{1}\big]
\le
\frac{L\big(f(\mathbf{x}_{0})-f^{\star}\big)}{p^{2}}\cdot\frac{1}{\sqrt{T}}
+\frac{p}{2L\sqrt{T}}
= \mathcal{O}\!\left(T^{-1/2}\right).
\]

Thus the average $\ell_{1}$‑norm of the true gradient decays at the rate $T^{-1/2}$.

%%%%%%%%%%%%%%%%%%%%%%%%%%%%%%%%%%%%%%%%%%%%%%%%%%%%%%%%%%%%
\section*{Special Cases}
\begin{enumerate}[label=\alph*.]
\item \textbf{Deterministic gradients ($\sigma=0$, $p=1$).}  
The bound reduces to the classic gradient descent guarantee
\[
\frac{1}{T}\sum_{k=0}^{T-1}\|\nabla f(\mathbf{x}_{k})\|_{1}
\le \frac{f(\mathbf{x}_{0})-f^{\star}}{\alpha T}
+\frac{L\alpha}{2},
\]
optimised by $\alpha=1/\sqrt{T}$ to obtain $\mathcal{O}(1/\sqrt{T})$.
\item \textbf{No momentum ($\beta=0$).}  
The theorem holds unchanged; the proof simply omits the term $-\alpha\beta\langle\nabla f,\mathbf{m}_{k}\rangle$.
\item \textbf{Large batch (variance $\sigma\to0$).}  
The correlation $p$ approaches $1$, allowing a larger admissible stepsize and a smaller constant factor in the rate.
\end{enumerate}

%%%%%%%%%%%%%%%%%%%%%%%%%%%%%%%%%%%%%%%%%%%%%%%%%%%%%%%%%%%%
\section*{Discussion}
The analysis shows that even when only the sign of a stochastic gradient is retained, the presence of an exponential moving average (momentum) yields a provable decrease of the objective under standard smoothness assumptions.  The convergence rate matches that of vanilla SGD up to a constant that depends on the sign‑correlation parameter $p$.  Consequently, Sign‑Momentum is theoretically justified for non‑convex smooth optimisation, while offering practical benefits such as reduced communication bandwidth in distributed settings.

\end{document}